\section{Conclusion and Future Work}\label{sec:conclusion}
This study has established that the reliability of few-shot anomaly thresholds depends on both the amount and the statistical unit of normal evidence. At nominal $\alpha=0.20$, target-only LOIO has reached an empirical FAR of 0.341 on Gaussian-corrupted MVTec at $k=4$, showing that an attainable operating point need not remain stable under shift. The finite floor $1/(k+1)$ has explained when a target-only rule is structurally silent, but crossing that floor has not supplied an exchangeability guarantee.

The category-count analysis has quantified the separate obstacle faced by source-assisted certification. Even after all-zero observed category losses, optimistic multiplicity-free distribution-free 95\% certification has required at least 14, 29, and 59 independent categories at $\alpha=0.20$, 0.10, and 0.05. The frozen allocation and Hoeffding rule have required more, whereas the strict audit has provided only three or four certification categories. Every category-certified target cell has consequently received the fail-closed threshold $\tau^\star=0$ across all 960 frozen configurations, whose shared boundary condition precludes interpreting them as independent trials. In contrast, image-unit analyses under the selected-mixture iid model have selected positive thresholds in 36.7\% to 60.3\% of target cells, but have targeted the selected source mixture rather than marginal risk over a new-category draw.

The practical contribution has therefore been a study-design discipline: distinguish ranking from threshold reliability, report attainable rather than nominal operating points, audit normal false alarms under declared shifts, define the certification unit before observing results, and count independent categories when the claim concerns a new category. Future work should collect category-rich source archives or introduce explicit parametric, hierarchical, or side-information assumptions with independently verified validity. It should also test real production shifts, temporal dependence, condition-estimation error, and end-to-end deployment cost. These directions can improve usefulness, but they should not relabel repeated images as category evidence, confuse marginal with category-conditional risk, or convert a source-domain certificate into unconditional target control.
